% =====================================================================
% DISPLAY ITEMS — NMI allows up to SIX in the main text for an
% Article. VERBATIM, Content Types: "Display items - up to 6 items
% (figures and/or tables)." An earlier version of this comment said eight
% and then ten; both were wrong. Extended Data is capped at TEN
% separately (Preparing your material: "A maximum of ten Extended Data
% display items is permitted"), and Supplementary Information is the
% uncapped tier.
%
% SO TWO OF THE EIGHT BELOW HAVE TO GO OR MERGE. Reviewers read Extended
% Data less carefully than the main text, so the six main slots are the
% real budget and are allocated here rather than discovered at the end.
%
% This file is the FIGURE LIST, not the figures. Each entry names what
% the panel shows, which arm it needs, and whether the script that makes
% it exists yet. The arm matters: a figure built from one arm and
% captioned as if it were the model is the error we just caught in the
% crop UMAPs.
% =====================================================================

\section*{Display items}

\note{BUDGET: SIX, not eight. Allocation below. Each figure must be regenerable
by one command from a committed script — that is reproducibility item R5
and it is cheaper to honour while building than to retrofit.}

\subsection*{Figure 1 — the confound, and the design}
\note{Schematic, not data. Why a per-donor offset and a batch effect are
the same geometric object, and where in the model the two are pulled
apart. The one figure a reader must understand to read any other. DOES
NOT EXIST; has to be drawn.}

\subsection*{Figure 2 — the frontier}
\note{The pre-registered factorial, the curve, and the two reported arms
marked on it. EXISTS (frontier plot, both arms).}

\subsection*{Figure 3 — \elaj\ against the field}
\note{Grouped bars, cross-donor agreement and the donor and study costs,
every panel oriented higher-is-better. EXISTS. Needs the width
limitation in the caption and \texttt{supervised\_on} populated for the
\elaj\ rows.}

\subsection*{Figure 4 — the corpus through the latent}
\note{The whole-corpus UMAP with its cell-type regions marked, as the
entry point to every crop that follows. EXISTS as
\texttt{01\_cell\_family}. ARM MUST BE NAMED IN THE CAPTION.}

\subsection*{Figure 5 — grouped by cell type, across tissue}
\note{The tissue crop. Stromal cells from five tissues together;
epithelium almost purely endometrial. EXISTS as \texttt{02\_tissue}.}

\subsection*{Figure 6 — the cycle is a cycle}
\note{THE HEADLINE FIGURE. Cycle-day crop on a cyclic colour map, with
the displacement numbers (1.07 / 0.98 / 0.39 / 0.34) beside it so the
visual and the measurement sit in one item. EXISTS as
\texttt{08\_cycleday}; the displacement panel does not and has to be
drawn.}

\subsection*{Figure 7 — disease signal, and where it is not}
\note{The immune result at 0.77/0.83 against 0.22 corpus-wide, AND the
duempelmann non-replication in the same figure. Deliberately one item,
because splitting them would let a reader take the positive without the
negative. PARTLY EXISTS (\texttt{03\_disease}).}

\subsection*{Figure 8 — ablations}
\note{C-6 first: the phase classifier removed. Reserved. DOES NOT EXIST
— the runs have not been done.}

% ---------------------------------------------------------------------
\section*{Extended Data}
\note{Unlimited, and where the honest overflow goes: the remaining crop
panels (phase, stage, grade, platform, progesterone), the per-arm
comparison tables, the sufficient-statistics tables at every coarsening,
the donor-overlap audit, and the baseline fitting details. Captions here
carry the detail that the main-text captions cannot afford.}

\section*{Supplementary Information}
\note{The accession table, the gene vocabulary and how it was derived,
the full factorial results, and the seeds. Things a reader consults
rather than reads.}
